\section{Related Work and Mappings: Details}
\label{app:related}

This appendix supplies the details behind the mappings of Section~\ref{sec:instances}.

\subsection{Invariant confluence (\S\ref{sec:iconfluence-instance})}
\label{app:iconfluence}

A \emph{transaction history} $h$ (here simply \emph{history}) is a finite
partial order of transaction events, each labeled by a transaction in $T$,
where $e'<e$ means that $e$ ran on a state that included the writes of $e'$.
A sub-history is \emph{closed} when it is downward closed: it contains every
event that precedes one of its own.  History
$h'$ \inrel $h$ when $h$ is a closed sub-history of $h'$.  Then the events
that $h'$ adds precede no event of $h$, and each event of $h$ keeps exactly
the predecessors it had.

The state of a history is computed event by event.  An event $e$, labeled by
the transaction $t_e$, runs on
$\mathrm{pre}(e)=s_0\cup\bigcup_{e'<e}\mathrm{post}(e')$ and produces
$\mathrm{post}(e)=t_e(\mathrm{pre}(e))$.  The state of $h$ is
$\mathrm{eval}(h)=s_0\cup\bigcup_{e\in h}\mathrm{post}(e)$.  Because $h$ is a
partial order and transactions are inflationary ($t(s)\supseteq s$),
$\mathrm{pre}(e)$ is the state of the sub-history strictly below $e$, and
$\mathrm{post}(e)$ is the state of the sub-history at or below $e$.  A
closed sub-history keeps each event's predecessors, so it also keeps each
event's $\mathrm{pre}$ and $\mathrm{post}$.  Hence the state of a union of
closed sub-histories is the union of their states.

A history is \emph{locally valid} when every $\mathrm{pre}(e)$ and
$\mathrm{post}(e)$ satisfies $\mathcal I$.  This is the case when replicas
abort transactions that would violate the invariant on their own state.  A
history is \emph{valid} when the state of each of its closed sub-histories
satisfies $\mathcal I$, that is, $\mathcal I(\mathrm{eval}(c))$ for every
closed sub-history $c$.

The spec $\Spec_{\mathcal I}$ has the locally valid histories as inputs;
a closed sub-history of a locally valid history is locally valid, so the
inputs are closed under the order.
An output is a state, and $u$ \outrel $s$ when $s\subseteq u$.  At input $h$,
each valid closed sub-history of $h$ contributes its state as a permitted
output:
\[
  P_{\mathcal I}(h)=\{\mathrm{eval}(g):g\text{ a valid closed
  sub-history of }h\}.
\]
Taking $g$ empty gives $s_0$, so the spec is total.

\begin{proof}[Proof of Proposition~\ref{prop:iconfluence-lm}]
We use the common-extension form of Lifted Monotonicity
(Proposition~\ref{prop:common-expansion-form}).  We first show, by induction
on $g$, that the state of every valid history $g$ is reachable.  Let $e$ be
maximal in $g$.  Then $g\setminus\{e\}$ and the sub-history strictly below $e$
are closed, hence valid.  By induction their states
$\mathrm{eval}(g\setminus\{e\})$ and $\mathrm{pre}(e)$ are reachable.  The
state $\mathrm{post}(e)$ is a valid transaction step from $\mathrm{pre}(e)$,
and $\mathrm{eval}(g)=\mathrm{eval}(g\setminus\{e\})\cup\mathrm{post}(e)$ is a
valid merge.

First assume that $T$ is invariant confluent.  Let rows
$(h_r,\mathrm{eval}(g_r))$, $r=1,\ldots,k$, have a common \indesc $j$.  Then each $g_r$ is a closed sub-history of $j$.  Let
$g=g_1\cup\cdots\cup g_k$; we show that $g$ is valid.  Every closed
sub-history $c$ of $g$ is the union of the pieces $c\cap g_r$.  Each piece is
a closed sub-history of the valid $g_r$, so its state is reachable.  Thus
$\mathrm{eval}(c)$ is a union of reachable states.
Invariant confluence makes each successive union valid, and a valid merge of
reachable states is reachable, so $\mathcal I(\mathrm{eval}(c))$.  Hence $g$
is valid, and $\mathrm{eval}(g)\in P_{\mathcal I}(j)$ \outrel every
$\mathrm{eval}(g_r)$.  For $k=0$, take $g$ empty.  Lifted Monotonicity
follows.

Now assume that $\Spec_{\mathcal I}$ satisfies Lifted Monotonicity.  Only its
directed fibers are used.  We first show that valid histories are closed under
disjoint union.  Let $g_1$ and $g_2$ be valid histories, and let $j$ be their
disjoint union, with events renamed apart and no order between the two
parts.  Each event keeps its
predecessors, so its $\mathrm{pre}$ and $\mathrm{post}$ are states of closed
sub-histories of a valid $g_r$.  So $j$ is locally valid, and
$\mathrm{eval}(g_1),\mathrm{eval}(g_2)\in P_{\mathcal I}(j)$.  Directed
fibers give a valid sub-history $g$ of $j$ with
$\mathrm{eval}(g)\supseteq\mathrm{eval}(g_1)\cup\mathrm{eval}(g_2)=\mathrm{eval}(j)\supseteq\mathrm{eval}(g)$.
So $\mathcal I(\mathrm{eval}(g_1)\cup\mathrm{eval}(g_2))$.  Every closed
$c\subseteq j$ is the disjoint union of $c\cap g_1$ and $c\cap g_2$, both
valid, and the same argument gives $\mathcal I(\mathrm{eval}(c))$.  So $j$ is
valid.  Valid histories are also closed under adding an event above all
others when its transaction step is valid.  By induction on derivations, every
reachable state is therefore the state of a valid history.  The pairwise
argument then gives $\mathcal I(s_1\cup s_2)$ for all reachable $s_1,s_2$.
Invariant confluence follows.
\end{proof}

\subsection{The HAT taxonomy (\S\ref{sec:hat-instance})}
\label{app:hat-taxonomy}

Bailis et al.\ classify transactional guarantees by what a system can provide
while every client that can reach a replica keeps getting
answers~\cite{bailis2013hat}.  The extended paper defines each guarantee by
the anomaly it rules out and gives algorithms for the guarantees that can be
provided~\cite{bailis2013hatextended}.  We use those definitions and
algorithms as written, except that where a definition disagrees with its own
example, we follow the example.

\paragraph{Inputs and outputs.}
The inputs extend the histories of Definition~\ref{def:hat-spec}.  Each item
has a fixed set of replicas, and a read is served by a replica of its item.  A
\emph{predicate read} reads a fixed range of items at once and returns a
version of each.  A transaction's reads are ordered.  A client runs one
transaction at a time.  Its transactions form its \emph{session}, and its
commit timestamps increase along the session.  Outputs are sets of read
answers, as before.

\paragraph{Guarantees.}
Each guarantee rules out one anomaly of the extended paper: a pattern in a
graph whose nodes are transactions.  The edges are the write-, read-, and
anti-dependencies of Adya on each item~\cite{adya1999weakconsistency}, and
session edges from each transaction to its client's later ones.  For example,
read your writes rules out a transaction that reads a version of $x$ older
than one that an earlier transaction of the same client wrote.  The extended
paper gives no formal definition of recency.  We take it to require that a
read return a version at least as new as every write by another transaction
whose commit happens before the read.  Because a replica can return only
versions it holds, the \emph{visibility store} permits an answer only if its
version is the initial one or arrived at the reading replica before the
read.

\begin{definition}[Guarantee spec, high availability, sticky read, sticky availability]
\label{def:hat-availability}
For a guarantee $G$, the table of $\Spec_G$ contains an output at an input when
the visibility store permits the output and the input extends to a complete
history that includes the output and has no anomaly that $G$ rules out.  The
guarantee $G$ is \emph{highly available} when some explicit relation
$E\subseteq P_G$ answers every read and some system for $\Spec_G$ has
Availability for $E$ and Partition Consistency.  A read is \emph{sticky} when
its replica has received every write of the client's earlier transactions to
the items that replica holds.  The guarantee $G$ is \emph{sticky available}
when the same holds for an $E$ that answers every sticky read.
\end{definition}

Sticky availability follows Bailis et al.: a client gets answers whenever its
transaction runs against a copy of the database that reflects all of its
earlier operations.  By Theorem~\ref{thm:complete-cap} and
Corollary~\ref{cor:complete-cap-any-system}, $G$ is highly available exactly
when $\Spec_G$ satisfies Lifted Monotonicity for some $E\subseteq P_G$ that
answers every read, and likewise for sticky availability.

\paragraph{Reveal rules.}
A transaction is \emph{pending stable} before a read when every replica of
every item it wrote has received its write before the read.  Bailis et al.'s
algorithm for monotonic atomic view reveals a write once its transaction is
pending stable.  For the session guarantees they also ``wait to reveal new
writes \ldots\ until each write's respective dependencies are visible on all
replicas.''  For cut isolation, a client caches the versions its transaction
has read.  The \emph{reveal rule} combines the three.  A write is revealed at
a read once its transaction and every transaction whose commit happens before
its commit are pending stable.  A transaction reads, for each item, the
highest version revealed at every one of its reads so far.  The \emph{sticky
reveal rule} answers only sticky reads, and it also reveals the client's own
earlier writes, which a sticky read's replica holds.

\begin{theorem}[HAT taxonomy]
\label{thm:hat-taxonomy}
\leavevmode
\begin{enumerate}
  \item Under the timestamp condition of Theorem~\ref{thm:hat-levels}, read
        committed and monotonic atomic view are highly available.  Item cut
        isolation, predicate cut isolation, monotonic reads, monotonic
        writes, and writes follow reads are highly available.
  \item Read your writes, PRAM, and causal consistency are sticky available
        and not highly available.  Given at least two participants and any
        assignment of clients to them, consider the deployment that lets each
        client's participant give the sticky reveal rule's answers to that
        client's reads.  Some system has Deployment Availability and Partition
        Consistency for it.
  \item Recency is neither highly available nor sticky available.  For any
        level that rules out Lost Update or Write Skew, no system has
        Availability for an explicit relation that answers every read and
        Partition Consistency.
\end{enumerate}
\end{theorem}

\begin{proof}[Proof sketch]
(1) The reveal rule answers every read.  Each answer depends only on the read's
past, so later input never changes it, and the rule is monotone.  A monotone
relation inside the spec satisfies Lifted Monotonicity for itself: rows below
$j$ have the relation's own answer at $j$ as a common \outnoun.  The rule
returns only versions the replica holds, and its answers have none of the
anomalies these guarantees rule out.  Consider monotonic reads.  If an earlier
transaction of a client read a version of $x$, that version was revealed.  It
had then reached every replica of $x$, so it stays revealed at every later
read.  The client's later transaction reads after the earlier one commits, so
it reads that version of $x$ or a newer one.
For read committed and monotonic atomic view, the timestamp condition rules
out G1c, as in Theorem~\ref{thm:hat-levels}.  By
Theorem~\ref{thm:complete-cap}, each guarantee is highly available.

(2) The sticky reveal rule does the same for every sticky read.  The deployment
authorizes only answers of this one monotone relation inside the spec, so it
is monotone across partitions, and Proposition~\ref{prop:deployment-complete-cap}
gives the system.  Without stickiness, consider a client that writes $x$ at
one replica and then reads $x$ at another replica that has not received the
write.  That replica holds only the initial version of $x$, which read your
writes rules out.  No explicit relation inside the spec answers the read.
PRAM and causal consistency include read your writes.

(3) For recency, another client's write of $x$ commits before a read but never
reaches the reading replica, which again holds only the initial version.  The
read is sticky, so stickiness does not help.  For Write Skew, two transactions
at partitioned replicas each read the item the other writes.  Each replica can
return only the initial version, and together the answers form the cycle that
Write Skew names, as for Lost Update in Theorem~\ref{thm:hat-levels}.
\end{proof}

\paragraph{The algorithm for monotonic atomic view.}
The extended paper also gives pseudocode for monotonic atomic
view~\cite[App.~B]{bailis2013hatextended}.  A replica keeps a write in a
\emph{pending} set until it has received acknowledgments from every replica of
the items its transaction wrote, and then moves it to a \emph{good} set.  A
read of $x$ may carry a required timestamp.  Without one, the replica returns
its highest good version of $x$.  With one, it returns a good version at least
as new, or else the pending write with exactly that timestamp.  After each
read, the client raises the required timestamp of every item that the
returned write's transaction also wrote, and it clears these timestamps when
the transaction ends.  The input does not fix which acknowledgments a
replica has received, so each replica's good set may be any set of pending
stable transactions that depends only on the read's past.  The pending case
returns only a write whose transaction an earlier read of the same transaction
found good, so every answer is a version the replica holds.

\begin{proposition}[The algorithm provides monotonic atomic view]
\label{prop:hat-mav-algorithm}
For every choice of the replicas' good sets, and under the timestamp condition
of Theorem~\ref{thm:hat-levels}, the answers of the algorithm for monotonic
atomic view form an explicit relation inside $P_{\mathrm{MAV}}$ that answers
every read and is monotone.  So monotonic atomic view is highly available.
\end{proposition}

\begin{proof}[Proof sketch]
Each answer depends only on the read's past, so the relation is monotone, and
it answers every read.  The timestamp condition rules out G1c, as for read
committed.  For the anomaly of monotonic atomic view, suppose a transaction
reads $x$ from $T_j$ and later reads $y$, which $T_j$ also wrote.  The first
read raised the required timestamp of $y$ to $T_j$'s, so the second read
returns $T_j$'s version of $y$ or a newer one.  Theorem~\ref{thm:complete-cap}
gives the system.
\end{proof}

\paragraph{Differences in form.}
Our classification matches theirs and differs in form in three ways.
First, our framework separates two reasons a guarantee can be unavailable.
The unavailable isolation levels fail Lifted Monotonicity: answers that the two
sides give have no common extension.  Read your writes without stickiness, and
recency in every case, fail earlier.  Some read has no permitted answer at all
(items (2) and (3) of the proof of Theorem~\ref{thm:hat-taxonomy}), so the
question of Lifted Monotonicity never arises.  Second, monotonic writes needs
the version order to follow each session's order.  Bailis et al.\ obtain this
from client timestamps, and our inputs assume it.  Third, we model a predicate
as a fixed range of items, and general predicates, whose items depend on the
data, are not modeled.  Bailis et al.\ also list several recency guarantees
without formal definitions, among them bounds on staleness and safe, regular,
and linearizable register semantics.  We model only the recency guarantee
defined above.

\subsection{Registers}
\label{app:registers}

Section~\ref{sec:complete-cap} derives the Gilbert--Lynch theorem from a
replica that is cut off from a completed write and must still answer a read.
The example here differs: two concurrent writes reach both replicas before the
partition.

\paragraph{The causal register.}
The causal register differs from the register of
Appendix~\ref{app:cap-response} in two ways.  Its input records causal
precedence instead of real-time precedence, and its output records the write
that supplies each read instead of the value returned.  An input is a finite operation history $h$: a set of reads and writes of one
register $x$ at replicas, with a strict partial order $\to$ of causal
precedence.  Causal precedence is program order at each replica, plus an edge
from each write to the operations a replica invokes after receiving it,
closed transitively.  A set of operations of $h'$ is \emph{causally closed}
when it contains every operation that causally precedes one of its members.
We set $h\InputOrd h'$ when the operations of $h$ form a causally closed
subset of $h'$ and $\to$ restricted to them is the order of $h$.  Thus an
\indesc never adds an operation to the causal past of an existing one.

An output is a \emph{source map} $\rho$: a partial map assigning to some reads
of $h$ a write of $h$ or the initial value $v_0$, ordered by extension.  A
\emph{serialization} for replica $p$ is a total order of $p$'s operations and
all writes of $h$ that extends $\to$.  It explains a read of $p$ with source
$w$ when $w$ is the last write before the read, or $v_0$ if there is none.
Under causal consistency~\cite{lamport1978time,ahamad1995causal},
$\rho\in\fiber{h}$ when each replica has a serialization explaining every read
of that replica in the domain of $\rho$.  Under linearizability, a single
total order of all operations of $h$ extending $\to$ must explain every read
in the domain~\cite{herlihy1990linearizability}.  Real-time edges beyond $\to$
play no role in the example below.

Two replicas start empty.  Participant $p$ completes $w_p(x,\dog)$ while $q$
concurrently completes $w_q(x,\catv)$.  The replicas exchange both writes.
Then a partition separates them, and each serves a local read, $r_p$ at $p$
and $r_q$ at $q$.  Both writes causally precede both reads, and the writes are
concurrent.

Under linearizability, a source map that sends $r_p$ to $w_q(x,\catv)$ needs a
linearization with $w_p(x,\dog)$ before $w_q(x,\catv)$.  One that sends $r_q$
to $w_p(x,\dog)$ needs the opposite order.  No source map \outrel both, so the
register is not monotone across the partition $\{\{p\},\{q\}\}$ under any
deployment authorizing these two reads.

Under causal consistency the two replicas may serialize the concurrent writes
differently, so one source map sends $r_p$ to $w_q(x,\catv)$ and $r_q$ to
$w_p(x,\dog)$.  More generally, consider the deployment in which each replica
answers only its own reads ($E_p$ holds the rows whose source map is defined
only on reads of $p$).  Under it the causal register is monotone across
partitions, for every history.  For each participant $p$ whose run emits, let
$\rho_p$ be the last source map it emits, at input $i_p$; then $\rho_p$
\outrel every output $p$ emitted.  Let $j$ be a common \indesc.  We extend a
serialization $L_p$ of $i_p$ explaining $\rho_p$ to $j$ by appending the new
operations of $p$ and the new writes in any order extending $\to$.  No new
operation causally precedes an old one, so the result extends $\to$ and
explains the same reads.  With these serializations for the replicas that
emitted, and arbitrary ones for the others, the union of the $\rho_p$ lies in
$\fiber{j}$ and \outrel every emitted output.
Separability is not used.  Two answers to the same read remain incompatible,
but only the reading replica may emit either.

\subsection{Real-time causal consistency}
\label{app:rtc}

Mahajan et al.\ showed that a system can provide real-time causal consistency,
but no stronger consistency, while staying always available and one-way
convergent~\cite{mahajan2011consistency}.  We recover both parts with real-time
precedence recorded in the input.  The argument rests on a property of
products of independent views.

\paragraph{Independent products.}
In the real-time causal register below, each replica independently chooses
a view of its own reads.  The facts we need hold for any such product of
independent views.  Fix participants
$\mathcal N$ with $|\mathcal N|\ge2$.  For each $p$ let $(W_p,\le_p)$ be a poset
and, for each input $i$, let $V_p(i)\subseteq W_p$ be a nonempty antichain, the
\emph{views} $p$ may expose at $i$.  A \emph{partial tuple} $x$ assigns to each
$p$ a value $x_p$ in $W_p$ or $\bot$, ordered coordinate-wise with $\bot$ below
every value; $x$ is a \emph{restriction} of a tuple $y$ when each coordinate
$x_p$ is $\bot$ or $y_p$.  The \emph{product interface} has a row
$(i,x)$ for every restriction $x$ of every $y\in\prod_pV_p(i)$.  A
\emph{strengthening} $T$ selects at each input a set of tuples in
$\prod_pV_p(i)$, and its rows are their restrictions.  It is \emph{strict} when
it omits a row of the product interface, and \emph{locally conservative} when
every $v\in V_p(i)$ is the $p$-coordinate of a tuple it selects at $i$.  Under
the \emph{owner deployment}, $p$ may emit exactly the rows whose tuple is
$\bot$ outside coordinate $p$.

\begin{theorem}[Independent products]
\label{thm:product-frontier}
\leavevmode
\begin{enumerate}
  \item If views continue (for $i\InputOrd j$ and $v\in V_p(i)$ some $w\in
        V_p(j)$ has $v\le_p w$), the product interface is monotone across
        partitions under the owner deployment.
  \item No strict, locally conservative strengthening is monotone across
        partitions under its owner deployment.
\end{enumerate}
\end{theorem}

\begin{proof}
(1) Fix a participant run family and a common \indesc $j$.  If $p$ belongs to
the family, the remainder table of $p$'s run has a row $(k_p,u)$, and $u$ \outrel every output $p$ emitted
(Lemma~\ref{lem:run-residual}).  As $u$ is a restriction of some
$y\in\prod_qV_q(k_p)$, continuation gives $w_p\in V_p(j)$ with $y_p\le_p w_p$.
Otherwise pick any $w_p\in V_p(j)$.  Each output of $p$ is $\bot$ outside
coordinate $p$, so $w$ \outrel every emitted output.
(2) Strictness gives $y\in\prod_pV_p(i)$ that $T$ does not select at $i$.  By
local conservativity, the tuple $e_p$ that is $y_p$ at $p$ and $\bot$
elsewhere is a row of $T$ at $i$, authorized for $p$.  Let every participant
emit $e_p$ at $i$; every partition separates two of them.  A common \outnoun
at $i$ is a restriction of a selected $z$ with $y_p\le_p z_p$ for every $p$, so
$y=z$ by the antichain property, a contradiction.
\end{proof}

\paragraph{Real-time causal views.}
Consider the histories of Appendix~\ref{app:registers} with a second strict
partial order $\prec$: $a\prec b$ when $a$ ended before $b$ began.  We assume
that $a\to b$ implies $a\prec b$, so replicas propagate only completed writes.
We also assume that $\to$ totally orders the operations of each replica.  One
input \inrel another when it does so as a history and $\prec$ agrees on the
old operations.  A write $w$ is \emph{visible} to a read $r$ when $w\to r$.  A
\emph{real-time causal view} of replica $p$ assigns to every read $r$ of $p$ a
source $s(r)$ such that
\begin{enumerate}
  \item $s(r)$ is a write visible to $r$, or $v_0$ if none is;
  \item $s(r)\not\prec w$ for every write $w$ visible to $r$;
  \item if $r'\to r$ is a read of $p$ to which $s(r)$ is visible, then
        $s(r')=s(r)$.
\end{enumerate}
Intuitively, these items are Mahajan et al.'s checks for the order that $p$'s
reads expose, in which each visible write precedes the read and each
overwritten write precedes the source.  Item~(1) says a replica returns only
writes it has received.  Item~(2) is their real-time check ``time doesn't
travel backward.''  Item~(3) keeps one replica's overwrites consistent.  Their
reads return all concurrent latest writes, while ours return one, so a view
matches their check when that write is the only latest one.

Views exist.  As each read's source we can take the visible write that comes
latest in a linear extension of $\prec$.  Views also continue.  An \indesc
changes neither a read's causal past nor $\prec$ on it, so an old read keeps
its source.  A new read takes as source its latest visible write in a linear
extension of $\prec$ together with the edges $x\prec s(r)$, one for each old
read $r$ and each write $x\ne s(r)$ visible to $r$.  All views of $p$ at an
input are defined on the same reads.  So no view extends a different one, and
the views form an antichain.  In the real-time causal product interface, each
replica answers its own reads.  By Theorem~\ref{thm:product-frontier}(1), this
interface is monotone across partitions under the owner deployment.

\paragraph{Maximality.}
A strengthening is \emph{real-time blind} when a tuple it selects at an input
stays selected at the same history with fewer real-time edges, provided its
views are still real-time causal there.  Mahajan et al.'s implementations do
not observe real time, so they return the same answers at both inputs.

\begin{proposition}[Real-time causal maximality]
\label{prop:rtc-maximal}
No strict, real-time blind strengthening of the real-time causal product
interface is monotone across partitions under its owner deployment.
\end{proposition}

\begin{proof}
By Theorem~\ref{thm:product-frontier}(2) it suffices to show local
conservativity.  Let $v$ be a view of $p$ at $(h,\prec)$.  We add real-time
edges that leave $v$ as $p$'s only view: for each read $r$ of $p$ and each
write $x\ne s(r)$ visible to $r$, we place $x$ before $s(r)$.  Formally, let
$a\prec_v b$ when $a\prec b$, or when $a=x$ or $a\prec x$, and $s(r)=b$ or
$s(r)\prec b$, for such $r$ and $x$.  Items~(2) and~(3) and program order make
$\prec_v$ a strict partial order, and it contains $\to$.  At $(h,\prec_v)$,
$v$ is $p$'s only view.  By~(1) a source is a visible write, and a visible
$x\ne s(r)$ has $x\prec_v s(r)$, which~(2) forbids.  Every input has a row of
the strengthening, and rows are restrictions of selected tuples.  So the
strengthening selects some tuple at $(h,\prec_v)$, and that tuple's
$p$-coordinate is $v$.  Real-time edges only exclude views, so the tuple's
views are real-time causal at $(h,\prec)$.  Blindness then selects the tuple
there.
\end{proof}

Each hypothesis has a counterpart in Mahajan et al.'s theorem.  Their always
availability corresponds to two conditions here: $T$ selects a tuple at every
input, so every replica can answer every read, and the system has Deployment
Availability.  One-way convergence needs no separate hypothesis: the input
records which writes each replica has received, the environment supplies that
input, and views must respect $\to$.  Real-time blindness is their assumption
that an implementation does not observe real time, which their proof uses to
swap the real-time order of concurrent writes.  The proof of
Proposition~\ref{prop:rtc-maximal} needs it too, since a strengthening may read
real time from the input, for instance to break ties between writes that
overlap in real time.
